% page 509 of the facsimile (image p-508 of the scan)
% block 165.1 x 263.6 mm ; left margin 36.9 mm
\pgc{15.240mm}{0.000mm}{
\l{}
\l{}
\l{}
\l{}
\l{\cel{52}501}
\l{}
\l{\sou{robinio}. (\sou{bot.}) Genero de arbori ek la familio \textquotedbl{}leguminosi\textquotedbl{}, di}
\l{\cel{3}qui la speco maxim konocata esas granda arboro kun flori}
\l{\cel{3}blanka, odoroza, qui kreskas grape, e di qua la ligno densa}
\l{\cel{3}esas apta por facar chari. - L.\cel{1}\sou{acacia}. -}
\l{}
\l{\sou{robro.}\cel{1}(en\cel{1}\sou{la wisto-ludo}). Partio ligita. -\cel{2}DEFR.}
\l{}
\l{\sou{robusta}. Di qua la konstituciono esas tre forta, tre vigoroza,}
\l{\cel{3}tre rezistiva. - DEFIS.}
\l{}
\l{\sou{rocheto}. (\sou{liturgio katol}.) Surpliso kun maniki rekta, ek stofo}
\l{\cel{3}blanka e tre dina, quan +weras la episkopi ed ula alt-ofi-}
\l{\cel{3}cieri di la eklezio katolika. - DEFIS.}
\l{}
\l{\sou{rodar}. (\sou{trans.})\cel{2}I. Entamar per mikra dento-mordi, beko-frapi. -}
\l{\cel{3}II. Konsumar, destruktar (korpo) per efiko lenta e nepercept-}
\l{\cel{3}eblaEFIS.}
\l{}
\l{\sou{rodio}. (\sou{kemio)}. Metalo harda e facile ruptebla, qua aspektas}
\l{\cel{3}quale aluminio, esas duktila e maleebla kande inkandecanta}
\l{\cel{3}til redeso ; lu fuzas ye 2000o C. -\cel{1}\sou{Simbolo kem}. Rh.}
\l{}
\l{\sou{rododendro}. (\sou{bot.}) Genero de planti ek la familio \textquotedbl{}erikacei\textquotedbl{},}
\l{\cel{3}arboreto sempre verda, remarkinda pro la beleso di lua flori.}
\l{\cel{3}- L.\cel{1}\sou{rhododendron}. - DEFIS.}
\l{}
\l{\sou{rodomont}o.\cel{2}Falsa bravo qua asertas ocidir multa enemiki ed}
\l{\cel{3}esar kapabla trafendar irgu. - EFI.}
\l{}
\l{\sou{rogo}. Amaso de ligno brancha sur qua la antiqui kustumis kom-}
\l{\cel{3}bustar sua mortinti. - IL.}
\l{}
\l{\sou{rogacioni}. (\sou{religio katol.}) Prego publika, agata procesione,}
\l{\cel{3}dum la tri dii qui preiras la Acenso-festo, por obtenar bona}
\l{\cel{3}rekolti. -}
\l{}
\l{\sou{roko}. I. Blokego de petro tre rezistiva qua retenesas da sulo. -}
\l{\cel{3}II. Blokego de petro harda, eskarpa. - III. (\sou{metaf.}) Quannulo}
\l{\cel{3}povas emocigar. - IV. (\sou{anat.)}\cel{2}Un ek la tri parti di la osto}
\l{\cel{3}temporala, tale nomizita pro lua hardeso. - EFIS.}
\l{}
\l{\sou{rokambolo}. (\sou{bot.}) Alio reda di Hispania, planto liliacea, perena,}
\l{\cel{3}bulboza. Lua bulbi saporas kelke quale alio, ma plu feble,}
\l{\cel{3}ed uzesas kom kondimento. - L.\cel{1}\sou{alliu}m sativum var. ophios-}
\l{\cel{3}\sou{corodon.}\cel{1}-\cel{2}DEF.}
\l{}
\l{\sou{rokoko.}\cel{2}I. Stilo ornementala qua esis segun-moda en la yarcento}
\l{\cel{3}18-esma. - II.(\sou{metaf.}) Stilo ornementala obsoleta, antiquatra.}
\l{\cel{3}- DEF.}
\l{}
\l{\sou{rolo.}\cel{1}I. (\sou{teatro)}. (a) Transskriburo di to quon aktoro devas}
\l{\cel{3}recitar ek teatrajo.\cel{2}- (b) To quon la aktoro devas recitar.}
\l{\cel{3}(c) Persono quan reprezentas la aktoro. - II. (\sou{metaf.)}}
\l{\cel{3}(a) Lo agenda kom funciono en ofico. - (b) Konduto-maniero}
\l{\cel{3}di ulu en ta o ca cirkonstanco. - DEFR.}
}
